\documentclass[a4paper,12pt,fleqn]{article}
\usepackage[ngerman]{babel}
\usepackage[utf8]{inputenc}
\usepackage[T1]{fontenc}
\usepackage{amsmath,amssymb}

\newcommand{\methodbox}[1]{%
\medskip
\noindent\fbox{%
    \begin{minipage}{0.94\linewidth}
    \textbf{Methodik / Definition.}\par\smallskip
    \small
    #1
    \end{minipage}%
}
\medskip\par
}

\newcommand{\examplebox}[1]{%
\medskip
\noindent\fbox{%
    \begin{minipage}{0.94\linewidth}
    \textbf{Beispiele.}\par\smallskip
    \small
    #1
    \end{minipage}%
}
\medskip\par
}

\newcommand{\taskbox}[1]{%
\medskip
\noindent\fbox{%
    \begin{minipage}{0.94\linewidth}
    \textbf{Aufgabenstellung.}\par\smallskip
    \small
    #1
    \end{minipage}%
}
\medskip\par
}

\begin{document}
\raggedright
\setlength{\parindent}{0pt}
\setlength{\mathindent}{0pt}

{\LARGE\bfseries ETI -- Aufgabenblatt 4 \\ Praesenzaufgaben\par}
\vspace{2em}

\section*{Aufgabe 4.1}

\taskbox{
Geben Sie eine CFG fuer die Sprache
\[
    \{a^n b^m c^m d^n\mid n,m\ge 0\}
\]
an.
}

\methodbox{
Die aeusseren Symbole $a$ und $d$ muessen paarweise gleich oft erzeugt werden.
Unabhaengig davon werden innen gleich viele $b$ und $c$ erzeugt. Deshalb trennen wir
die beiden Zaehler in zwei Nonterminale.
}

Eine passende kontextfreie Grammatik ist
\[
G=(\{S,T\},\{a,b,c,d\},R,S)
\]
mit den Produktionen
\[
\begin{aligned}
S &\to aSd \mid T,\\
T &\to bTc \mid \varepsilon.
\end{aligned}
\]

Nach $n$ Anwendungen von $S\to aSd$ entsteht aussen $a^n$ und $d^n$.
Danach erzeugt $T$ durch $m$ Anwendungen von $T\to bTc$ genau $b^m c^m$ und beendet
mit $T\to\varepsilon$. Somit gilt
\[
    L(G)=\{a^n b^m c^m d^n\mid n,m\ge 0\}.
\]

\section*{Aufgabe 4.2}

\taskbox{
Betrachten Sie
\[
T=\{0^i b\mid i\ge 1,\ b\in\{0,1\}^*,\ |b|_0\le i\}
\]
und die CFG $G=(V,\Sigma,R,S)$ mit
\[
V=\{S,X,Y\},\qquad \Sigma=\{0,1\}
\]
und Regeln
\[
\begin{aligned}
S&\to XY,\\
X&\to 0X\mid 0,\\
Y&\to Y1\mid 1Y\mid Y0\mid 0Y\mid 1\mid 0\mid \varepsilon.
\end{aligned}
\]
Bestimmen Sie $L(G)$ und vergleichen Sie diese Sprache mit $T$.
}

Das Nonterminal $X$ erzeugt genau eine nichtleere Folge von Nullen:
\[
    L(X)=\{0^i\mid i\ge 1\}=0^+.
\]

Das Nonterminal $Y$ erzeugt alle Woerter ueber $\{0,1\}$. Einerseits sind
$0,1,\varepsilon$ direkt erzeugbar. Andererseits duerfen beliebig oft links oder rechts
eine $0$ oder eine $1$ angefuegt werden. Also gilt
\[
    L(Y)=\{0,1\}^*.
\]

Damit ist
\[
    \boxed{L(G)=0^+\{0,1\}^*}.
\]

\medskip
Jedes Wort aus $T$ beginnt mit mindestens einer $0$. Deshalb gilt
\[
    T\subseteq L(G).
\]
Die Inklusion ist echt. Zum Beispiel liegt
\[
    0100\in L(G),
\]
denn es beginnt mit einer $0$. Fuer $T$ ist aber nur die Zerlegung
$0100=0^1\cdot 100$ moeglich, weil direkt danach eine $1$ kommt. Dabei enthaelt
$100$ zwei Nullen, also
\[
    |100|_0=2>1.
\]
Folglich ist $0100\notin T$ und damit
\[
    \boxed{T\subsetneq L(G)}.
\]

\section*{Aufgabe 4.3}

\taskbox{
Diskutieren und begruenden Sie, ob die Sprachen aus der vorherigen Aufgabe
kontextfrei oder sogar regulaer sind.
}

\methodbox{
Regulaere Sprachen sind unter Konkatenation, Vereinigung und Schnitt mit regulaeren
Sprachen abgeschlossen. Kontextfreie Sprachen koennen durch CFGs beschrieben werden.
Um Nicht-Regulaerheit zu zeigen, schneiden wir mit einer einfachen regulaeren Sprache
und wenden das Pumping-Lemma an.
}

\textbf{Die Sprache $L(G)$.}

Aus Aufgabe 4.2 wissen wir
\[
    L(G)=0^+\{0,1\}^*.
\]
Diese Sprache ist regulaer, zum Beispiel beschrieben durch den regulaeren Ausdruck
\[
    00^*(0\cup 1)^*.
\]
Damit ist $L(G)$ insbesondere auch kontextfrei.

\medskip
\textbf{Die Sprache $T$.}

Die Sprache $T$ ist kontextfrei, denn Aufgabe 4.4 gibt explizit eine CFG fuer $T$ an.
Sie ist jedoch nicht regulaer.

\medskip
Angenommen, $T$ waere regulaer. Dann waere auch der Schnitt mit der regulaeren Sprache
$0^*10^*$ regulaer:
\[
K:=T\cap 0^*10^*
  =\{0^i10^j\mid i\ge 1,\ 0\le j\le i\}.
\]

Sei $p$ die Pumping-Laenge von $K$. Betrachte
\[
    w=0^p10^p\in K.
\]
Nach dem Pumping-Lemma gibt es eine Zerlegung $w=xyz$ mit
$|xy|\le p$, $|y|\ge 1$ und $xy^rz\in K$ fuer alle $r\ge 0$.
Wegen $|xy|\le p$ besteht $y$ nur aus Nullen aus dem ersten Block, also
$y=0^\ell$ mit $\ell\ge 1$.

Pumpe nun nach unten, also $r=0$:
\[
    xz=0^{p-\ell}10^p.
\]
Hier stehen nach der $1$ mehr Nullen als davor. Also ist
\[
    p>p-\ell,
\]
und damit $xz\notin K$. Das widerspricht dem Pumping-Lemma. Also ist $K$ nicht
regulaer und folglich kann $T$ nicht regulaer sein.

\section*{Aufgabe 4.4}

\taskbox{
Geben Sie eine Grammatik $H$ mit $L(H)=T$ an.
}

\methodbox{
Jede fuehrende $0$ erzeugt ein Guthaben fuer hoechstens eine $0$ im Restwort $b$.
Ein Hilfsnonterminal erzeugt daher entweder nur Einsen oder genau eine Null mit
beliebig vielen Einsen davor und danach.
}

Wir definieren
\[
H=(\{S,C,D\},\{0,1\},R,S)
\]
mit Produktionen
\[
\begin{aligned}
S &\to 0SC \mid 0C,\\
C &\to 1C \mid 0D \mid \varepsilon,\\
D &\to 1D \mid \varepsilon.
\end{aligned}
\]

Dabei gilt
\[
    L(C)=1^*\ \cup\ 1^*01^*.
\]
Also erzeugt jedes $C$ hoechstens eine Null. Das Startsymbol $S$ erzeugt bei jeder
Rekursion zuerst eine fuehrende $0$ und haengt spaeter ein $C$ an. Nach $i\ge 1$
fuehrenden Nullen entstehen genau $i$ solche $C$-Bloecke:
\[
    S\Rightarrow^* 0^i C_i C_{i-1}\cdots C_1.
\]
Da jeder $C$-Block hoechstens eine Null erzeugt, enthaelt der erzeugte Rest hoechstens
$i$ Nullen. Somit ist jedes erzeugte Wort in $T$.

\medskip
Umgekehrt sei $w\in T$. Dann gibt es eine Zerlegung
\[
    w=0^i b
\]
mit $i\ge 1$ und $|b|_0\le i$. Zerlege $b$ entlang seiner Nullen in hoechstens $i$
Teile, sodass jeder Teil entweder nur aus Einsen besteht oder genau eine Null enthaelt.
Jeder dieser Teile wird von $C$ erzeugt; fehlende Teile erzeugen $\varepsilon$.
Damit kann $H$ genau dieses Wort $w$ ableiten.

Folglich gilt
\[
    \boxed{L(H)=T}.
\]

\end{document}
